\section{Grenzen}

\begin{itemize}
  \item \textbf{Nur 700 Zeilen.} Die Kennzahlen tragen Bootstrap-Intervalle
        bzw. Fold-Streuung, sind aber von begrenzter Praezision.
  \item \textbf{Kein Zeitbezug.} Es gibt keinen Beobachtungszeitpunkt, zu
        dem die Merkmale erhoben wurden. Das Modell sagt eine Zugehoerigkeit
        vorher, kein Ereignis zu einem Zeitpunkt.
  \item \textbf{Die 450 ungelabelten Zeilen bleiben ungenutzt.} Sie tragen
        keine Zielvariable; ein halbueberwachter Ansatz waere moeglich, aber
        nur mit zusaetzlicher Annahme.
  \item \textbf{Das Kostenmodell.} Das Verhaeltnis 50:1 ist die vom
        Auftraggeber genannte Relation, keine Messung aus den Daten. Die
        Szenarien 1:1 bis 10:1 zeigen, wie empfindlich das Optimum auf diese
        Annahme reagiert.
  \item \textbf{Keine Ursachenaussage.} Die Koeffizienten zeigen
        Zusammenhang, nicht Wirkung.
\end{itemize}

\section{Reproduktion}

\begin{lstlisting}[language=bash]
python3 data/fetch_data.py       # Rohdaten holen + Form pruefen
python3 src/01_describe.py       # Kennzahlen + Abbildungen
python3 src/02_models.py         # Modellvergleich (verschachtelte CV)
python3 src/03_drivers.py        # Treiberanalyse
python3 src/04_threshold.py      # Schwellenwert + Kostenmodell
python3 src/05_significance.py   # gepaarter Wilcoxon-Test
python3 src/report.py            # dieses PDF (LaTeX ueber Tectonic)
\end{lstlisting}

Die Ergebnisdateien liegen in \texttt{results/} (JSON und CSV im deutschen
Power-BI-Format), die Abbildungen in \texttt{figures/}.
